%=====================================================================
\section{Appendix: Model Code}\label{sec:code}
%=====================================================================
This appendix lists the code of the model exactly as it appears in the notebook \texttt{gcn\_elliptic.ipynb}
(cells that only draw figures or print statistics are left out). Everything is plain PyTorch; no pre-built graph layer is used.

\subsection*{A.1 \ Setup and imports}
Libraries, random seed and plot colours.
\begin{lstlisting}
import os, time, random
import numpy as np
import pandas as pd
import scipy.sparse as sp
from scipy.sparse.csgraph import connected_components
import torch
import torch.nn as nn
import torch.nn.functional as F
import matplotlib.pyplot as plt
import networkx as nx
from sklearn.linear_model import LogisticRegression
from sklearn.ensemble import RandomForestClassifier
from sklearn.metrics import (precision_score, recall_score, f1_score, accuracy_score,
                             confusion_matrix, average_precision_score)

def set_seed(seed):
    random.seed(seed); np.random.seed(seed); torch.manual_seed(seed)

set_seed(42)
torch.set_num_threads(os.cpu_count())
os.makedirs("figures", exist_ok=True)

# colours used in every figure
C_LICIT, C_ILLICIT, C_UNKNOWN = "#2a78d6", "#eb6834", "#b4b2ab"
C_MODEL = {"GCN": "#2a78d6", "Skip-GCN": "#eb6834", "MLP": "#1baf7a", "Random Forest": "#4a3aa7",
           "Logistic Regression": "#e87ba4"}
plt.rcParams.update({"figure.dpi": 110, "axes.spines.top": False, "axes.spines.right": False,
                     "axes.grid": True, "grid.alpha": 0.25, "font.size": 10})

print("PyTorch", torch.__version__, "| NumPy", np.__version__, "| pandas", pd.__version__)
\end{lstlisting}

\subsection*{A.2 \ Loading the Elliptic dataset}
Reads the three CSV files (features, labels, edges). If they are missing, PyTorch Geometric is used only to download them.
\begin{lstlisting}
FILES = ["elliptic_txs_features.csv", "elliptic_txs_classes.csv", "elliptic_txs_edgelist.csv"]

def locate_raw_files(data_dir="data/elliptic_bitcoin_dataset"):
    if all(os.path.exists(os.path.join(data_dir, f)) for f in FILES):
        return data_dir
    from torch_geometric.datasets import EllipticBitcoinDataset   # downloader only
    EllipticBitcoinDataset("data/pyg_elliptic")
    return "data/pyg_elliptic/raw"

raw_dir = locate_raw_files()
t0 = time.time()
feat_df  = pd.read_csv(os.path.join(raw_dir, FILES[0]), header=None)
class_df = pd.read_csv(os.path.join(raw_dir, FILES[1]))
edge_df  = pd.read_csv(os.path.join(raw_dir, FILES[2]))
print(f"loaded from {raw_dir} in {time.time()-t0:.1f}s")
print("features:", feat_df.shape, "| classes:", class_df.shape, "| edges:", edge_df.shape)
\end{lstlisting}

\subsection*{A.3 \ Features, labels and edges}
Turns the CSV files into a feature matrix, a label vector ($-1$ unknown, 0 licit, 1 illicit) and the list of money flows.
\begin{lstlisting}
tx_ids    = feat_df[0].values
time_step = feat_df[1].values.astype(int)
X_all     = feat_df.iloc[:, 2:].values.astype(np.float32)        # 165 features
LOCAL, AGG = slice(0, 93), slice(93, 165)                         # 93 local + 72 aggregated

N = len(tx_ids)
pos = pd.Series(np.arange(N), index=tx_ids)                        # txId -> row index

label_str = class_df.set_index("txId").loc[tx_ids, "class"].astype(str).values
y_np = np.full(N, -1, dtype=np.int64)          # -1 unknown, 0 licit, 1 illicit
y_np[label_str == "2"] = 0
y_np[label_str == "1"] = 1

src = pos[edge_df["txId1"].values].values
dst = pos[edge_df["txId2"].values].values

n_ill, n_lic, n_unk = (y_np == 1).sum(), (y_np == 0).sum(), (y_np == -1).sum()
print(f"transactions (nodes): {N:,}")
print(f"directed edges       : {len(src):,}")
print(f"features per node    : {X_all.shape[1]} (93 local + 72 aggregated)")
print(f"time steps           : {time_step.min()} .. {time_step.max()}")
print(f"illicit / licit / unknown : {n_ill:,} / {n_lic:,} / {n_unk:,}")
print(f"illicit share of labelled : {n_ill/(n_ill+n_lic):.2%}")
\end{lstlisting}

\subsection*{A.4 \ Train/test split and feature scaling}
Train on time steps 1--34, test on 35--49. Features are scaled with the training-period mean and standard deviation.
\begin{lstlisting}
train_mask_np = (y_np >= 0) & (time_step <= 34)
test_mask_np  = (y_np >= 0) & (time_step >= 35)

mu = X_all[time_step <= 34].mean(0); sd = X_all[time_step <= 34].std(0) + 1e-6
X_std = (X_all - mu) / sd

X_af = torch.tensor(X_std)                 # all 165 features
X_lf = torch.tensor(X_std[:, LOCAL])       # 93 local features only
y    = torch.tensor(y_np)
train_mask, test_mask = torch.tensor(train_mask_np), torch.tensor(test_mask_np)

for name, m in [("train (t 1-34)", train_mask_np), ("test (t 35-49)", test_mask_np)]:
    print(f"{name}: {m.sum():6,} labelled | illicit {(y_np[m]==1).sum():5,} ({(y_np[m]==1).mean():.1%})")
print(f"\nAccuracy of always predicting 'licit' on test: {(y_np[test_mask_np]==0).mean():.2%}  (illicit F1 = 0)")
\end{lstlisting}

\subsection*{A.5 \ Building the normalised adjacency matrix $\hat A$}
Adds self-loops, makes the graph undirected, and applies $\hat A = \tilde D^{-1/2}\tilde A\tilde D^{-1/2}$ as a sparse matrix. The identity version turns the GCN into the MLP baseline.
\begin{lstlisting}
def build_norm_adj(src, dst, n):
    src, dst = torch.as_tensor(src), torch.as_tensor(dst)
    loops = torch.arange(n)
    row = torch.cat([src, dst, loops])          # both directions + self-loops  -> A~ = A + I
    col = torch.cat([dst, src, loops])
    key = torch.unique(row * n + col)           # remove duplicate edges
    row, col = key // n, key % n
    d = torch.zeros(n).scatter_add_(0, row, torch.ones(row.numel()))   # D~_ii
    w = d[row].pow(-0.5) * d[col].pow(-0.5)                             # D~^-1/2 A~ D~^-1/2
    return torch.sparse_coo_tensor(torch.stack([row, col]), w, (n, n)).coalesce()

def identity_adj(n):                            # A_hat = I  turns the GCN into an MLP
    i = torch.arange(n)
    return torch.sparse_coo_tensor(torch.stack([i, i]), torch.ones(n), (n, n)).coalesce()

A_hat = build_norm_adj(src, dst, N)
I_hat = identity_adj(N)
print("A_hat:", tuple(A_hat.shape), "| non-zeros:", f"{A_hat._nnz():,}",
      f"| density {A_hat._nnz()/N**2:.2e}")
print("dense float32 size would be", f"{N*N*4/1e9:.0f} GB")

# sanity check: A_hat is symmetric and its largest eigenvalue is 1
v = torch.sparse.sum(A_hat, 1).to_dense()
print("row sums: min %.3f  max %.3f" % (v.min(), v.max()))
\end{lstlisting}

\subsection*{A.6 \ The GCN model}
Each layer computes $\hat A(HW) + b$. The same class builds the plain GCN, the Skip-GCN and deeper GCNs.
\begin{lstlisting}
class GCNLayer(nn.Module):
    '''H_out = A_hat @ (H W) + b'''
    def __init__(self, in_dim, out_dim):
        super().__init__()
        self.W = nn.Parameter(torch.empty(in_dim, out_dim))
        self.b = nn.Parameter(torch.zeros(out_dim))
        nn.init.xavier_uniform_(self.W)
    def forward(self, A, H, AH=None):
        if AH is not None:                       # A_hat @ H already known:  A_hat (H W) = (A_hat H) W
            return AH @ self.W + self.b
        return torch.sparse.mm(A, H @ self.W) + self.b

class GCN(nn.Module):
    def __init__(self, in_dim, hid_dim, out_dim=2, n_layers=2, dropout=0.0, skip=False):
        super().__init__()
        dims = [in_dim] + [hid_dim] * (n_layers - 1) + [out_dim]
        self.layers = nn.ModuleList(GCNLayer(dims[k], dims[k + 1]) for k in range(n_layers))
        self.skip = nn.Linear(in_dim, out_dim, bias=False) if skip else None
        self.dropout = dropout
    def forward(self, A, X, AX=None, return_hidden=False):
        H, hidden = X, None
        for k, layer in enumerate(self.layers):
            if k == 0 and AX is not None and self.dropout == 0:
                H = layer(A, X, AH=AX)           # first layer re-uses the pre-computed A_hat X
            else:
                H = layer(A, F.dropout(H, self.dropout, self.training))
            if k < len(self.layers) - 1:
                H = F.relu(H); hidden = H
        if self.skip is not None:
            H = H + self.skip(X)
        return (H, hidden) if return_hidden else H

m = GCN(165, 100)
print(m)
print("trainable parameters:", sum(p.numel() for p in m.parameters()))
\end{lstlisting}

\subsection*{A.7 \ Loss, training loop and metrics}
Weighted cross-entropy (0.3 licit / 0.7 illicit), Adam optimiser, and precision / recall / F1 / AP on the test period.
\begin{lstlisting}
CLASS_W = torch.tensor([0.3, 0.7])
HP = dict(hid=100, lr=1e-2, epochs=200, wd=5e-4, dropout=0.0)

def scores(y_true, y_pred, p_illicit=None):
    out = dict(precision=precision_score(y_true, y_pred, zero_division=0),
               recall=recall_score(y_true, y_pred, zero_division=0),
               f1=f1_score(y_true, y_pred, zero_division=0),
               micro_f1=accuracy_score(y_true, y_pred))
    if p_illicit is not None:
        out["ap"] = average_precision_score(y_true, p_illicit)
    return out

def train_gcn(A, X, seed=42, n_layers=2, skip=False, log_every=0, record=False, **hp):
    hp = {**HP, **hp}
    set_seed(seed)
    model = GCN(X.shape[1], hp["hid"], 2, n_layers=n_layers, dropout=hp["dropout"], skip=skip)
    opt = torch.optim.Adam(model.parameters(), lr=hp["lr"], weight_decay=hp["wd"])
    AX = torch.sparse.mm(A, X)       # A_hat and X never change, so A_hat X is computed once
    hist = []
    ytr, yte = y[train_mask], y[test_mask]
    for ep in range(1, hp["epochs"] + 1):
        model.train()
        out = model(A, X, AX)
        loss = F.cross_entropy(out[train_mask], ytr, weight=CLASS_W)
        opt.zero_grad(); loss.backward(); opt.step()
        if record and (ep == 1 or ep % 10 == 0):
            model.eval()
            with torch.no_grad():
                pred = model(A, X, AX).argmax(1)
            hist.append(dict(epoch=ep, loss=loss.item(),
                             train_f1=f1_score(ytr, pred[train_mask]),
                             test_f1=f1_score(yte, pred[test_mask]),
                             train_acc=(pred[train_mask] == ytr).float().mean().item(),
                             test_acc=(pred[test_mask] == yte).float().mean().item()))
            if log_every and (ep == 1 or ep % log_every == 0):
                h = hist[-1]
                print(f"epoch {ep:4d} | loss {h['loss']:.4f} | train F1 {h['train_f1']:.3f} | "
                      f"test F1 {h['test_f1']:.3f} | test acc {h['test_acc']:.3f}")
    model.eval()
    with torch.no_grad():
        logits, hidden = model(A, X, AX, return_hidden=True)
        prob = logits.softmax(1)[:, 1]
    return model, prob.numpy(), hidden, pd.DataFrame(hist)

def eval_probs(prob, mask=test_mask_np):
    return scores(y_np[mask], (prob[mask] >= 0.5).astype(int), prob[mask])
\end{lstlisting}

\subsection*{A.8 \ Training the GCN}
Trains the 2-layer GCN for 200 epochs with seed 42 and prints the test scores.
\begin{lstlisting}
t0 = time.time()
gcn, p_gcn, H_gcn, hist_gcn = train_gcn(A_hat, X_af, seed=42, record=True, log_every=100)
print(f"\ntraining time: {time.time()-t0:.0f}s")
print("initial loss:", round(hist_gcn.loss.iloc[0], 4), " (ln 2 = 0.6931 for a 50/50 guess)")
res_gcn = eval_probs(p_gcn)
print({k: round(v, 4) for k, v in res_gcn.items()})
\end{lstlisting}

\subsection*{A.9 \ Baselines and repeated runs}
Trains Logistic Regression, Random Forest, MLP, GCN and Skip-GCN on all features and on local features, with 5 seeds each.
\begin{lstlisting}
def run_sklearn(kind, X, seed=42):
    if kind == "Logistic Regression":
        clf = LogisticRegression(max_iter=2000)
    else:
        clf = RandomForestClassifier(n_estimators=50, max_features=50, random_state=seed, n_jobs=-1)
    clf.fit(X[train_mask_np], y_np[train_mask_np])
    prob = np.zeros(N); prob[test_mask_np] = clf.predict_proba(X[test_mask_np])[:, 1]
    return clf, prob

SEEDS = [42, 1, 2, 3, 4]
FEATS = {"AF": (X_std, X_af), "LF": (X_std[:, LOCAL], X_lf)}
NEURAL = [("MLP", I_hat, False), ("GCN", A_hat, False), ("Skip-GCN", A_hat, True)]

rows, probs = [], {}
t0 = time.time()
for feats, (Xn, Xt) in FEATS.items():
    _, p = run_sklearn("Logistic Regression", Xn)
    probs[("Logistic Regression", feats)] = p
    rows.append(dict(model="Logistic Regression", features=feats, seed=0, **eval_probs(p)))
    for s in SEEDS:
        _, p = run_sklearn("Random Forest", Xn, seed=s)
        if s == 42: probs[("Random Forest", feats)] = p
        rows.append(dict(model="Random Forest", features=feats, seed=s, **eval_probs(p)))
        for kind, A, skip in NEURAL:
            if (kind, feats, s) == ("GCN", "AF", 42):
                p = p_gcn                                    # already trained in Section 7
            else:
                _, p, _, _ = train_gcn(A, Xt, seed=s, skip=skip)
            if s == 42: probs[(kind, feats)] = p
            rows.append(dict(model=kind, features=feats, seed=s, **eval_probs(p)))
        print(f"{feats} seed {s:2d} done ({time.time()-t0:.0f}s)")

multi_df = pd.DataFrame(rows)
multi_df.to_csv("figures/all_runs.csv", index=False)
\end{lstlisting}
